\title{Controlling Text-to-Image Diffusion by Orthogonal Finetuning}

\begin{abstract}
Large-scale text-to-image diffusion models have demonstrated remarkable performance in creating lifelike images from textual descriptions. However, efficiently steering or customizing these robust models for a range of downstream tasks remains a significant and unresolved issue. To address this, we present a theoretically grounded adaptation technique termed Orthogonal Finetuning (OFT) for modifying text-to-image diffusion models in new application contexts. In contrast to prior approaches, OFT can formally maintain hyperspherical energy, a property representing the pairwise correlations among neurons constrained to the unit hypersphere. Our investigations reveal that this energy preservation is vital to maintain the semantic fidelity of generative outputs in text-to-image diffusion systems. To further enhance the robustness of the finetuning process, we propose Constrained Orthogonal Finetuning (COFT), which introduces an explicit constraint on the hypersphere’s radius. We focus on two major text-to-image adaptation tasks: subject-driven generation, which entails synthesizing images of a given entity using a limited set of reference images and textual descriptions, and controllable generation, which requires the model to respond to auxiliary conditioning inputs. Empirical results indicate that our OFT methodology surpasses established techniques regarding both image quality and training efficiency.
\end{abstract}

\section{Introduction}

State-of-the-art text-to-image diffusion models~\cite{saharia2022photorealistic,ramesh2022hierarchical,rombach2022high} exhibit excellent abilities in text-based, high-resolution image creation. Nevertheless, although textual input offers strong generative guidance, it often falls short in achieving precise, fine-level control of visual content. To mitigate this limitation, we concentrate on two specific categories of text-to-image generation tasks in this work:

\begin{itemize}[leftmargin=*,nosep]

    \item \textbf{Subject-driven generation}~\cite{ruiz2023dreambooth}: This setting involves generating new images featuring a specified subject in varied contexts, conditioned on a handful of images depicting that subject and an accompanying text prompt.
    \item \textbf{Controllable generation}~\cite{zhang2023adding,mou2023t2i}: In this scenario, the model generates images that conform both to a textual description and to an external control signal (such as canny edge maps or segmentation layouts).
\end{itemize}

The central challenge in both tasks involves how to adapt text-to-image diffusion models through finetuning while maintaining their generative capabilities learned during pretraining. An effective finetuning strategy should meet the following criteria: (1) \emph{training efficiency}, which refers to minimizing the number of trainable parameters and the required training epochs, and (2) \emph{generalizability preservation}, meaning the model should retain its ability to generate high-fidelity and diverse outputs. In pursuit of these goals, finetuning is most commonly performed by either incrementally updating the existing neuron weights with a low learning rate (\emph{e.g.}, \cite{ruiz2023dreambooth}) or by introducing a lightweight module with re-parameterized neuron weights (\emph{e.g.}, \cite{hulora2022,zhang2023adding}). However, these straightforward approaches are insufficient for reliably preserving the generative power attained in pretraining. There also remains a gap in the theoretical basis for designing optimal finetuning schemes or determining suitable hyperparameters, such as the number of epochs. This difficulty largely arises because there is no concrete metric for evaluating how well the finetuned model retains its pretrained generative strengths. Current finetuning practices often rest on the implicit assumption that smaller Euclidean distances between the parameters of the finetuned and pretrained models correspond to better preservation of pretrained abilities, which motivates the common use of very small learning rates. Although this may sometimes be valid, we contend that relying solely on Euclidean distance fails to adequately capture the semantic aspects of model preservation. Consequently, adopting a more structurally informed measure to assess the discrepancy between the finetuned and pretrained models could substantially enhance both the retention of pretrained generative performance and the stability of the finetuning process.

Motivated by empirical findings that hyperspherical similarity effectively captures semantic relationships~\cite{liu2017deep,liu2018decoupled,chen2020angular}, we adopt the notion of hyperspherical energy~\cite{liu2018learning} to describe the pairwise dependencies among neurons within a layer. Hyperspherical energy is computed as the aggregate of hyperspherical similarity measures (\emph{e.g.}, cosine similarity) for all neuron pairs in a layer, thus quantifying how uniformly the neurons are distributed over the unit hypersphere~\cite{liu2021learning}. We posit that, for effective finetuning, the hyperspherical energy of the model should closely match that of the original pretrained version. A straightforward strategy might introduce a regularization term to enforce consistency of hyperspherical energy during finetuning; however, this approach does not inherently guarantee that the energy difference will be minimized. Instead, we exploit a fundamental property of hyperspherical energy—namely, that applying the same orthogonal transformation to all neurons preserves pairwise hyperspherical similarities. Building upon this invariance, we introduce Orthogonal Finetuning~(OFT), a method designed to adapt large text-to-image diffusion models for specific downstream tasks while maintaining constant hyperspherical energy. The core principle involves optimizing a shared orthogonal transformation per layer so that pairwise neuron angles remain unchanged. This approach can be interpreted as redefining the coordinate system in which neurons within a layer are represented. Recognizing that a lower Euclidean distance between the finetuned and pretrained models further enhances the retention of pretrained capabilities, we also propose an extension—Constrained Orthogonal Finetuning~(COFT)—that restricts the finetuned solution to lie within a fixed-radius hypersphere centered around the pretrained neuron configurations.

The underlying rationale for utilizing orthogonal transformation in neuron finetuning draws partial inspiration from the 2D Fourier transform, which separates an image into its magnitude and phase spectra. Notably, the phase spectrum—characterizing the angular relation between the input and basis—retains most of the image’s semantic content. For instance, combining the phase spectrum from an original image with an arbitrary magnitude spectrum is often sufficient to reconstruct the image’s semantics. This observation indicates that semantic manipulation of a generated image—central to both subject-driven and controllable generation—is fundamentally achieved by adjusting neuron directions. Nevertheless, altering neuron orientations without constraint risks degrading the generative performance established during pretraining. To address this, we advocate for maintaining the angular relationships between neuron pairs, guided by the hypothesis that these inter-neuron angles encode essential network knowledge. Consequently, it is intuitive to learn a layer-wise, shared orthogonal transformation for neurons, ensuring that the hyperspherical energy remains consistent.

Additionally, our approach is informed by orthogonal over-parameterized training~\cite{liu2021orthogonal}, which develops classification networks by orthogonally transforming randomly initialized networks. The motivation lies in the property that such networks typically have low expected hyperspherical energy, and the objective of \cite{liu2021orthogonal} is to preserve low hyperspherical energy throughout training—a factor linked to improved generalization in classification~\cite{liu2018learning,lin2020regularizing}. The results in \cite{liu2021orthogonal} demonstrate that orthogonal transformations possess enough capacity to optimize classification networks effectively. However, our work addresses a distinct setting: the finetuning of text-to-image diffusion models, targeting enhanced controllability and superior downstream generative capabilities. There are two principal differences between OFT and the framework of \cite{liu2021orthogonal}. Firstly, whereas \cite{liu2021orthogonal} prioritizes minimizing hyperspherical energy, OFT is designed to conserve the hyperspherical energy of the pretrained model to preserve its internal representations during finetuning. Actively reducing this energy in diffusion model finetuning risks damaging original semantic structures. Secondly, OFT explicitly minimizes deviation from the pretrained network, motivating a constrained formulation; in contrast, \cite{liu2021orthogonal} introduces no such restrictions. Successful finetuning hinges on balancing flexibility with the preservation of the model’s foundational structure, and we posit that our OFT framework strikes this balance effectively. The main contributions of our work are summarized as follows:

\begin{itemize}[leftmargin=*,nosep]

    \item We introduce a new finetuning strategy, Orthogonal Finetuning (OFT), designed to enhance the controllability of text-to-image diffusion models. For greater finetuning stability, we present a constrained extension that restricts the angular displacement relative to the original pretrained parameters.
    \item In contrast to prior finetuning approaches, OFT is distinctive in its ability to update model parameters while analytically maintaining hyperspherical energy, a property we empirically demonstrate as being critical to upholding the pretrained model’s semantic generative integrity.
    \item We utilize OFT for two distinct applications: subject-driven image generation and controllable synthesis. Through extensive experimental evaluations, we showcase that OFT yields clear gains over state-of-the-art methods with respect to generation fidelity, convergence rate, and stability throughout finetuning. OFT also exhibits enhanced sample efficiency, requiring far fewer training instances and epochs to reach strong convergence.
    \item For the task of controllable generation, we propose a novel control consistency metric that quantitatively assesses controllability. This method involves inferring the control signal from synthesized images and comparing it to the original control directive, thereby further corroborating OFT’s superior controllability.
\end{itemize}

\section{Related Work}

\textbf{Text-to-image diffusion models}. Recent advancements in text-to-image synthesis~\cite{saharia2022photorealistic,ramesh2022hierarchical,rombach2022high,gu2022vector,nichol2022glide} are primarily attributable to innovations in diffusion-based generative models~\cite{ho2020denoising,song2021score,song2021denoising,dhariwal2021diffusion} and joint vision-language modeling~\cite{su2020vlbert,lu2019vilbert,radford2021learning,wang2021simvlm,li2022blip,li2023blip,alayrac2022flamingo,schuhmann2022laion}. Models such as GLIDE~\cite{nichol2022glide} and Imagen~\cite{saharia2022photorealistic} carry out diffusion in pixel space; GLIDE couples training of its text encoder and diffusion prior on paired data, whereas Imagen employs a fixed, pretrained text encoder throughout. Approaches like Stable Diffusion~\cite{rombach2022high} and DALL-E2~\cite{ramesh2022hierarchical} instead utilize latent spaces for their diffusion processes: Stable Diffusion leverages VQ-GAN~\cite{esser2021taming} to form a codebook-based latent space, and DALL-E2 incorporates CLIP~\cite{radford2021learning} to derive a unified image-text latent embedding. Beyond diffusion frameworks, the field also encompasses generative adversarial networks~\cite{reed2016generative,zhang2017stackgan,xu2018attngan,li2019controllable} and autoregressive techniques~\cite{ramesh2021zero,ding2021cogview,wu2022nuwa,yuscaling} for text-conditional image generation. OFT, being model-agnostic, is compatible with any diffusion-based text-to-image generative architecture.

\textbf{Subject-driven generation}. In attempts to limit subject alterations, techniques such as~\cite{avrahami2022blended,nichol2022glide} incorporate user-supplied masks as additional conditioning input. Inversion methods~\cite{choi2021ilvr,dhariwal2021diffusion,ramesh2022hierarchical,gal2022image} focus on modifying backgrounds or contextual elements while retaining the subject. Alternatively, \cite{hertz2022prompt} achieve both localized and global manipulations without the need for explicit masks. However, these strategies often fail to accurately conserve the subject’s identity-defining features. Pivotal Tuning~\cite{roich2022pivotal} mitigates this by finetuning the generator near an inverted code and imposing a regularization term to better maintain subject identity. Likewise, \cite{nitzan2022mystyle} establish a subject-specific generative prior from an individual's face image collection. Instance-level variation is also pursued in~\cite{casanova2021instance}, though it may come at the expense of identity-specific fidelity. DreamBooth~\cite{ruiz2023dreambooth} customizes the diffusion model by integrating a unique token and optimizing reconstruction and prior-preservation losses with only a small set of subject images. Building on the DreamBooth setup, OFT replaces the conventional finetuning—typically done with reduced learning rates—with parameter updates governed by orthogonal transformations.

\textbf{Controllable generation}. Image-to-image translation represents a subtype of controllable generation, where most prior approaches employ conditional generative adversarial networks~\cite{isola2017image,zhu2017toward,wang2018high,park2019semantic,choi2018stargan}. Recent works have also leveraged diffusion models in this context~\cite{saharia2022palette,voynov2022sketch,wang2022pretraining}. Among newer advances, ControlNet~\cite{zhang2023adding} enables precise control over a pretrained diffusion model by finetuning it with auxiliary control signals, resulting in strong performance for guided generation. Similarly, the concurrent T2I-Adapter~\cite{mou2023t2i} technique enhances image controllability through the finetuning of a pretrained diffusion model. Building upon these works~\cite{zhang2023adding,mou2023t2i}, we employ OFT to finetune pretrained diffusion models, and observe improved controllability while requiring less annotated data and fewer tunable parameters. Importantly, OFT maintains test-time computational efficiency by not adding any extra inference costs.

\textbf{Model finetuning}. The adaptation of large pretrained models to downstream applications has seen substantial uptake in recent years~\cite{devlin2018bert,brown2020language,he2022masked}. Adaptation techniques—including those in~\cite{hulora2022,houlsby2019parameter,pfeiffer2020adapterfusion}—are especially prominent in natural language processing research. LoRA~\cite{hulora2022}, which bears the closest resemblance to OFT, introduces a low-rank additive update for model weights during finetuning. By contrast, OFT applies layer-wise orthogonal transformations to the neurons in a multiplicative fashion, which, by construction, keeps the pairwise angular relationships between neurons within a layer unchanged—leading to greater finetuning stability.

\section{Orthogonal Finetuning}

\subsection{Why Does Orthogonal Transformation Make Sense?}

We begin by elucidating the rationale for employing orthogonal transformations in the finetuning of text-to-image diffusion models. We parse this overarching question into two components: (1) the motivation for modifying neuron directions (\emph{i.e.}, their angles), and (2) the suitability of using orthogonal transformations to enact these changes.

Addressing the first question, we are motivated by empirical findings in \cite{liu2018decoupled,chen2020angular} indicating that the angular difference in features effectively reflects the semantic disparity. The SphereNet~\cite{liu2017deep} work further demonstrates that constraining neural activations to the unit hypersphere results in comparable model expressiveness and, in some cases, improved generalization, suggesting that the directionality of neuron vectors efficiently encapsulates the principal data representations. To further illustrate the significance of neuron angle, we design a simple experiment depicted in Figure~\ref{fig:toy_ae}, wherein a standard convolutional autoencoder is trained on a dataset of flower images. During training, the conventional inner product is employed for feature map computation ($z$ denotes the convolution output given kernel $\bm{w}$ and sliding window input $\bm{x}$). In the test phase, three distinct approaches to feature map generation are contrasted: (a) replicating the inner product from training, (b) considering only magnitude, and (c) extracting angular information. The outcomes, as presented in Figure~\ref{fig:toy_ae}, reveal that reconstructing input images using neuron angular information nearly achieves perfect recovery, whereas magnitude alone provides negligible information. Notably, cosine-based activation (c) is not used in training; all training relies solely on the inner product. These observations suggest that neuron directions are chiefly responsible for encoding the semantic content of input data. Thus, adjusting neuron directions appears to be a promising strategy for effectively altering semantic attributes of images.

For the second question, the most direct method for adjusting neuron direction involves synchronously applying a rotation or reflection to all neurons within the same layer, inherently leveraging an orthogonal transformation. Although more adaptable angular transformations—allowing varied rotations across neurons—could be devised, our investigations indicate that orthogonal transformations achieve an optimal balance of adaptability and structural regularization. In support, \cite{liu2021orthogonal} establishes the robust learning potential provided by orthogonal transformations in neural networks. To empirically underscore our viewpoint, we implement a controllable generation task within the ControlNet~\cite{zhang2023adding} framework; results are illustrated in Figure~\ref{fig:ortho}. The experimental findings indicate that OFT, under the orthogonality constraint, consistently produces stable and accurate outcomes after twenty training epochs, whereas omitting the constraint results in failed generation and no meaningful control. These results underscore the critical role of the orthogonality constraint within OFT.

\subsection{General Framework}

The standard approach to finetuning employs gradient descent, typically with a reduced learning rate, to modify the parameters of a model or certain submodules. This use of a small learning rate inherently restricts how much the finetuned model $\bm{M}$ can diverge from its original pretrained configuration $\bm{M}^0$, leading to an implicit minimization of $\thickmuskip=2mu \medmuskip=2mu \| \bm{M} - \bm{M}^0\|$. While this implicit regularization aligns with intuition, it may allow too much flexibility, especially when adapting large models. LoRA addresses this by incorporating an explicit low-rank limitation on the weight changes, requiring that $\thickmuskip=2mu \medmuskip=2mu \text{rank}(\bm{M} - \bm{M}^0)=r'$, with $r'$ set to a small constant. OFT, on the other hand, enforces a constraint on the similarity among neuron pairs: $\thickmuskip=2mu \medmuskip=2mu \| \text{HE}(\bm{M})-\text{HE}(\bm{M}^0) \|=0$, where $\text{HE}(\cdot)$ calculates the hyperspherical energy of a weight matrix. 

For illustration, let us consider a dense layer $\thickmuskip=2mu \medmuskip=2mu \bm{W}=\{\bm{w}_1,\cdots,\bm{w}_n\}\in\mathbb{R}^{d\times n}$, where each $\bm{w}_i\in\mathbb{R}^{d}$ denotes an individual neuron, and let $\bm{W}^0$ be the pretrained weights. The output $\thickmuskip=2mu \medmuskip=2mu \bm{z}\in\mathbb{R}^n$ is computed as $\thickmuskip=2mu \medmuskip=2mu \bm{z}=\bm{W}^\top\bm{x}$, for input vector $\thickmuskip=2mu \medmuskip=2mu \bm{x}\in\mathbb{R}^d$. Within this context, OFT can be seen as attempting to reduce the discrepancy in hyperspherical energy between the updated and initial weights:

\begin{equation}\label{eq:oft_principle}
\footnotesize
\min_{\bm{W}} \left\| \text{HE}(\bm{W})-\text{HE}(\bm{W}^0)\right\|~~\Leftrightarrow~~\min_{\bm{W}} \bigg{\|} \sum_{i\neq j} \|\hat{\bm{w}}_i-\hat{\bm{w}}_j\|^{-1}-\sum_{i\neq j} \|\hat{\bm{w}}^0_i-\hat{\bm{w}}^0_j\|^{-1}\bigg{\|}
\end{equation}

Here, $\thickmuskip=2mu \medmuskip=2mu \hat{\bm{w}}_i:=\bm{w}_i/\|\bm{w}_i\|$ refers to the normalized $i$-th neuron, and the hyperspherical energy $\text{HE}(\bm{W})$ of the layer $\bm{W}$ is given by $\thickmuskip=2mu \medmuskip=2mu \text{HE}(\bm{W}):= \sum_{i\neq j} \|\hat{\bm{w}}_i-\hat{\bm{w}}_j\|^{-1}$. It is evident that the lowest achievable value for Eq.~\eqref{eq:oft_principle} is zero. This minimum occurs whenever $\bm{W}$ differs from $\bm{W}^0$ solely by an orthogonal transformation; in other words, $\thickmuskip=2mu \medmuskip=2mu \bm{W}=\bm{R}\bm{W}^0$, where $\thickmuskip=2mu \medmuskip=2mu \bm{R}\in\mathbb{R}^{d\times d}$ is an orthogonal matrix—its determinant indicating whether the transformation is a rotation ($1$) or reflection ($-1$). OFT is precisely based on this insight: rather than modifying all neuron weights freely, finetuning is limited to learning shared orthogonal transformations for each layer. Form

\begin{equation}\label{eq:oft_general}
    \footnotesize
    \bm{z} = \bm{W}^\top \bm{x} = (\bm{R} \cdot \bm{W}^0)^\top \bm{x},~~~\text{s.t.}~\bm{R}^\top\bm{R} = \bm{R}\bm{R}^\top = \bm{I}
\end{equation}

Here, $\bm{W}$ is the weight matrix updated by OFT, while $\bm{I}$ refers to the identity matrix. An overview of OFT can be seen in Figure~\ref{fig:block}. In line with LoRA’s zero initialization, it is necessary for OFT to adjust the pretrained model strictly starting from $\bm{W}^0$. Thus, the orthogonal matrix $\bm{R}$ is initialized as an identity matrix, ensuring that the adapted model commences from the original pretrained parameters. To maintain orthogonality for $\bm{R}$ throughout training, differential orthogonalization approaches, as outlined in \cite{liu2021orthogonal,lezcano2019cheap}, are applicable. We will elaborate on methods to enforce orthogonality efficiently in computational terms.

\subsection{Efficient Orthogonal Parameterization}\label{sect:eff_ortho}

Conventional techniques for producing orthogonal matrices, such as the differentiable Gram-Schmidt process, are generally considered computationally intensive and not scalable in practice~\cite{liu2021orthogonal}. To overcome these computational challenges, we utilize Cayley parameterization, enabling the efficient construction of an orthogonal matrix. Concretely, we form the orthogonal matrix as $\thickmuskip=2mu \medmuskip=2mu \bm{R} = (\bm{I} + \bm{Q})(I - \bm{Q})^{-1}$, where $\bm{Q}$ is a skew-symmetric matrix satisfying $\thickmuskip=2mu \medmuskip=2mu \bm{Q} = -\bm{Q}^\top$. The increased efficiency does introduce a minor constraint—Cayley parameterization restricts $\bm{R}$ to special orthogonal matrices, i.e., those with determinant $1$. Our experiments indicate that this constraint has negligible impact on model performance. However, even when employing the Cayley transform for parameterizing $\bm{R}$, the total parameter count may remain excessive when $d$ is large. To mitigate this, we suggest using a block-diagonal configuration for $\bm{R}$, partitioning it into $r$ blocks, which yields the following expression:

\begin{equation}
    \footnotesize
    \bm{R} = \text{diag}(\bm{R}_1, \bm{R}_2, \cdots, \bm{R}_r) = \begin{bmatrix}
    \bm{R}_1 \in \text{O}(\frac{d}{r}) &  & \\
    & \ddots & \\
    & & \bm{R}_r \in \text{O}(\frac{d}{r})
    \end{bmatrix} \in \text{O}(d)
\end{equation}

Here, $\text{O}(d)$ represents the group of $d$-dimensional orthogonal matrices. We denote $\thickmuskip=2mu \medmuskip=2mu \bm{R}\in\mathbb{R}^{d\times d}$ as well as each block $\thickmuskip=2mu \medmuskip=2mu \bm{R}_i\in\mathbb{R}^{d/r\times d/r}$ for all $i$ as orthogonal matrices. When $\thickmuskip=2mu \medmuskip=2mu r=1$, the block-diagonal structure reduces to a standard, unconstrained orthogonal matrix. In general, a $d\times d$ orthogonal matrix is parameterized by $\thickmuskip=2mu \medmuskip=2mu d(d-1)/2$ independent variables, yielding a complexity of $\mathcal{O}(d^2)$. For a block-diagonal orthogonal matrix with $r$ blocks, the parameter count becomes $\thickmuskip=2mu \medmuskip=2mu d(d/r-1)/2$, so the overall complexity is reduced to $\thickmuskip=2mu \medmuskip=2mu \mathcal{O}(d^2/r)$. Further parameter reduction can be achieved by tying (sharing) blocks across the diagonal, specifically by letting $\thickmuskip=2mu \medmuskip=2mu\bm{R}_i=\bm{R}_j$ for all $i\neq j$, resulting in only $\mathcal{O}(d^2/r^2)$ parameters. Importantly, these modifications retain the orthogonality of $\bm{R}$, thereby maintaining the same level of hyperspherical energy.

Next, we analyze parameter efficiency by comparing OFT with LoRA. LoRA introduces $r'$ trainable parameters per direction, so its total number of trainable weights is $\thickmuskip=2mu \medmuskip=2mu r'(d+n)$. If both $r$ and $r'$ are proportional to the neuron count $d$, say $\thickmuskip=2mu \medmuskip=2mu r=r'=\alpha d$ for some constant $0<\alpha\leq 1$, then the parameter complexity for LoRA scales as $\mathcal{O}(d^2+dn)$, whereas for OFT, the complexity is reduced to $\mathcal{O}(d)$. To clarify this comparison with a practical scenario, consider a $128\times 128$ weight matrix: LoRA with $\thickmuskip=2mu \medmuskip=2mu r'=8$ has $2{,}048$ trainable parameters, while OFT with $\thickmuskip=2mu \medmuskip=2mu r=8$ (without shared blocks) requires $960$ parameters.

\subsection{Constrained Orthogonal Finetuning}

The expressive capacity of standard OFT may be further curtailed by restricting the adapted parameters to reside within a compact region around the original pretrained parameters. This leads to COFT, which carries out the forward computation as follows:

\begin{equation}\label{eq:coft_general}
    \footnotesize
    \bm{z}=\bm{W}^\top\bm{x}=(\bm{R}\cdot\bm{W}^0)^\top\bm{x},~~~\text{s.t.}~\bm{R}^\top\bm{R}=\bm{R}\bm{R}^\top=\bm{I},~\norm{\bm{R}-\bm{I}}\leq \epsilon
\end{equation}

This approach enforces both orthogonality and proximity (in terms of a deviation $\epsilon$) to the identity matrix. Orthogonality can be accomplished using the Cayley transform described in Section~\ref{sect:eff_ortho}. However, directly applying the $\epsilon$-proximity condition to Cayley-parameterized matrices introduces significant challenges. To facilitate understanding, we can examine the Cayley transform through the lens of a Neumann series, where we approximate $\thickmuskip=2mu \medmuskip=2mu \bm{R}=(\bm{I}+\bm{Q})(I-\bm{Q})^{-1}$ by $\thickmuskip=2mu \medmuskip=2mu \bm{R}\approx\bm{I}+2\bm{Q}+\mathcal{O}(\bm{Q}^2

Because the series converges in the operator norm, we can equivalently express the constraint $\thickmuskip=2mu \medmuskip=2mu\|\bm{R}-\bm{I}\|\leq\epsilon$ by pushing it through the Cayley transform. The resulting requirement becomes $\thickmuskip=2mu \medmuskip=2mu \|\bm{Q}-\bm{0}\|\leq\epsilon'$, where $\bm{0}$ designates a zero matrix and $\epsilon'$ is a separate error tolerance parameter. Projected gradient descent can conveniently enforce this revised constraint on $\bm{Q}$. For identity initialization of the orthogonal matrix $\bm{R}$, one simply initializes $\bm{Q}$ with zeros. COFT thus unifies two clear constraints: a minimal hyperspherical energy change, and a bounded deviation from the original, pretrained model. While most existing finetuning strategies implicitly employ the second constraint, COFT distinguishes itself by imposing it explicitly. Our experiments indicate that this explicit constraint enables COFT to achieve superior finetuning stability relative to OFT, in addition to its strong overall performance. Figure~\ref{fig:eps_coft} illustrates how varying $\epsilon$ alters COFT's performance. Larger $\epsilon$ values permit greater flexibility, with COFT's output approximating that of OFT; conversely, reducing $\epsilon$ draws the COFT-finetuned model closer in behavior to the original pretrained text-to-image diffusion model.

\subsection{Re-scaled Orthogonal Finetuning}

We introduce a straightforward modification to OFT by allowing each neuron to have its own learnable scaling parameter. This extension is motivated by the observation that scaling neuron outputs leaves the hyperspherical energy invariant due to subsequent normalization. The modified forward propagation is defined as: $\thickmuskip=2mu \medmuskip=2mu\bm{z}=(\bm{R}\bm{W}^0\bm{D})^\top\bm{x}$\footnote{\textbf{Errata}: In the NeurIPS camera ready submission, the formula for the forward computation of re-scaled OFT was incorrectly written as $\thickmuskip=2mu \medmuskip=2mu\bm{z}=(\bm{D}\bm{R}\bm{W}^0)^\top\bm{x}$. The implementation, however, is correct; thus, the results in Appendix~\ref{app:rescaled} remain valid.} where $\thickmuskip=2mu \medmuskip=2mu \bm{D}=\text{diag}(s_1,\cdots,s_n)\in\mathbb{R}^{n\times n}$ is a diagonal matrix with all entries $s_1,\cdots,s_n$ learnable and strictly positive. Unlike the original OFT in Eq.~\eqref{eq:oft_general}, where only $\bm{R}$ is optimized, this variant makes both the scaling matrix $\bm{D}$ and the orthogonal matrix $\bm{R}$ trainable. This re-scaled OFT variant further enhances the adaptability of OFT with only a marginal increase in parameters. For experimental clarity, we primarily report results for the unscaled OFT to highlight the role of orthogonal transformations alone, but empirically observe that the re-scaled variant often provides superior performance (see Appendix~\ref{app:rescaled}).

\section{Intriguing Insights and Discussions}

\textbf{OFT is architecture-independent}. In theory, OFT is applicable to any neural network architecture. For Transformer models, finetuning methods like LoRA are commonly used on the attention layers~\cite{hulora2022}. To maintain a fair comparison, our experiments confine OFT finetuning to the attention weights as well. OFT's structure also lends itself naturally to convolutional layer finetuning: the block-diagonal form of $\bm{R}$ yields interpretations specific to convolutional operations—unlike LoRA. For instance, selecting the number of blocks equal to the number of input channels results in a situation where each block acts on an individual neuron channel, effectively functioning like depth-wise convolution kernels~\cite{chollet2017xception}. If all blocks of $\bm{R}$ are shared, the orthogonal transformation is applied identically across all channels. Preliminary experiments with OFT on convolutional layers are discussed in Appendix~\ref{app:convolution}.

\textbf{Connection to LoRA}. LoRA limits the extent of changes to the pretrained weight matrix by incorporating a low-rank matrix, thereby maintaining stability in the pretrained information. In contrast, OFT preserves the integrity of pretrained weights by applying an orthogonal (and thus full-rank) transformation, ensuring that the core pretraining information is not lost. The forward computation in OFT can be recast as $\thickmuskip=2mu \medmuskip=2mu \bm{z}=(\bm{R}\bm{W}^0)^\top\bm{x}=(\bm{W}^0+(\bm{R}-\bm{I})\bm{W}^0)^\top\bm{x}$, where $\thickmuskip=2mu \medmuskip=2mu (\bm{R}-\bm{I})\bm{W}^0$ plays a role similar to the low-rank update mechanism in LoRA. When $\bm{W}^0$ is of full rank, OFT acts analogously to a low-rank update provided that $\thickmuskip=2mu \medmuskip=2mu \bm{R}-\bm{I}$ is itself low-rank. Analogous to the rank parameter $r'$ in LoRA, OFT introduces a block-diagonal parameter $r$ to constrain the count of learnable parameters. Notably, LoRA and OFT exemplify fundamentally different approaches to parameter efficiency: LoRA utilizes low-rankness to achieve compactness, while OFT leverages sparsity, specifically block-diagonal orthogonality, to limit parameter overhead.

\textbf{Why OFT converges faster}. As shown in Figure~\ref{fig:toy_ae}, altering the orientation of neurons most effectively influences semantic change, which is precisely the focus of OFT. Furthermore, OFT frames the finetuning process as optimization on a smooth hypersphere manifold, a perspective associated with favorable optimization landscapes. This benefit has also been substantiated empirically by \cite{liu2021orthogonal}.

\textbf{Why not minimize hyperspherical energy}. A significant distinction from \cite{liu2021orthogonal} lies in our decision not to pursue hyperspherical energy minimization. For classification tasks, one aims to avoid redundant neurons. Achieving minimal hyperspherical energy results in neurons being distributed evenly across the hypersphere, which is counterproductive during finetuning, as such uniformity can undermine valuable information acquired through pretraining.

\textbf{Balancing flexibility and regularization in finetuning}. Our analysis reveals a fundamental balance between model flexibility and regularization. While conventional finetuning offers the highest level of flexibility, it suffers from limited robustness and often leads to instability or collapse of the model. Notably, the simplicity of OFT allows it to achieve a well-calibrated compromise between flexibility and regularity by maintaining the angles between neuron pairs. The block-diagonal approach acts as a more stringent form of regularization on the orthogonal matrix.

\textbf{No increase in inference cost}. In contrast to ControlNet, our OFT framework maintains zero additional computational cost at inference for the finetuned model. At inference, the orthogonal matrix $\bm{R}$ learned during training is multiplied directly with the original weight matrix $\bm{W}^0$ to yield the final weights $\thickmuskip=2mu \medmuskip=2mu \bm{W}=\bm{R}\bm{W}^0$. Therefore, inference performance remains identical to that of the pretrained model.

\section{Experiments and Results}

\textbf{Experimental setup}. We utilize Stable Diffusion v1.5~\cite{rombach2022high} as the pretrained backbone for text-to-image synthesis. For rigorous evaluation, generated samples from each method are selected randomly. In subject-driven generation experiments, procedures align closely with those of DreamBooth~\cite{ruiz2023dreambooth}. For controlled generation tasks, we follow established setups from ControlNet~\cite{zhang2023adding} and T2I-Adapter~\cite{mou2023t2i}. To permit a fair comparison against LoRA, we restrict application of OFT or COFT exclusively to the same model layers targeted by LoRA. Additional experimental details and supplementary results can be found in Appendix~\ref{app:settings}.

\subsection{Subject-driven Generation}

\textbf{Setup}. We employ DreamBooth~\cite{ruiz2023dreambooth} and LoRA~\cite{hulora2022} as reference baselines. All compared approaches use the identical loss formulation as in DreamBooth. For DreamBooth and LoRA, experimental procedures and hyperparameters are selected based on their respective original works. Expanded results and further discussion are available in Appendix~\ref{app:settings},\ref{app:coft_oft},\ref{app:more_qualitative},\ref{app:failure}.

\textbf{Finetuning stability and convergence}. To assess the stability and speed of convergence during finetuning, we compare DreamBooth, LoRA, OFT, and COFT, with visual summaries provided in Figure~\ref{fig:teaser} and Figure~\ref{fig:db_convergence}. The empirical findings indicate that both COFT and OFT reliably finetune the diffusion model without exhibiting instability. At the 400-iteration mark, both DreamBooth and OFT-based methods exhibit strong controllability, whereas LoRA demonstrates poor subject identity retention. By 2000 iterations, DreamBooth tends toward image collapse, and LoRA is unable to accurately depict the yellow shirt, generating yellow fur instead. In contrast, OFT and COFT maintain robust and stable control over the output images throughout the process. These outcomes reinforce the efficacy of the OFT framework in providing rapid and stable convergence for subject-driven image generation. It is worth emphasizing that reduced sensitivity to the specific number of finetuning iterations is particularly advantageous, as it lessens the burden of iteration-tuning for varying subjects. With OFT and COFT, it is feasible to assign a generously high number of iterations from the outset without the need for detailed tuning. Moreover, with an appropriate choice of $\epsilon$ in COFT, setting both the learning rate and the number of iterations becomes a straightforward task.

\textbf{Quantitative comparison}. Adopting the evaluation protocols from \cite{ruiz2023dreambooth}, we quantitatively assess subject fidelity (using DINO~\cite{caron2021emerging} and CLIP-I~\cite{radford2021learning}), adherence to text prompts (CLIP-T~\cite{radford2021learning}), and image diversity (LPIPS~\cite{zhang2018unreasonable}). The CLIP-I metric is calculated as the mean pairwise cosine similarity between CLIP embeddings of generated images and their real counterparts. For DINO, we utilize ViT S/16 DINO features and employ an analogous procedure. CLIP-T measures the average cosine similarity between embeddings of generated images and their corresponding text prompts. LPIPS is used to compute the mean cosine similarity between generated outputs conditioned on the same prompt and subject. According to Table~\ref{table:dreambooth_final}, both COFT and OFT considerably surpass DreamBooth and LoRA in the DINO and CLIP-I subject fidelity assessments, and perform comparably or marginally better in text alignment and diversity evaluations. Each method undergoes repeated finetuning of the same pretrained model using 30 distinct random seeds to reduce randomness effects. Collectively, these results underscore that the OFT framework not only offers improved stability and convergence during training, but also secures superior and consistently reproducible performance outcomes.

\textbf{Qualitative comparison}. To provide a clearer and more intuitive illustration of OFT's strengths, we present several randomly selected subject-driven generation samples in Figure~\ref{fig:dreambooth_final}. For consistency and fairness, each example across different methods is generated using the same finetuned model, ensuring that no separate prompt engineering is performed for individual outputs. We pick, for every method, the model yielding the highest validation CLIP score. As displayed in Figure~\ref{fig:dreambooth_final}, both OFT and COFT are highly effective in preserving semantic information related to the subject, whereas LoRA frequently struggles to maintain the correct subject identity (\emph{e.g.}, the bowl sample in which LoRA fails to retain any aspects of the subject's identity). Furthermore, OFT and COFT demonstrate substantially stronger correspondence between the generated image and the textual prompt compared to DreamBooth, which, while typically able to keep the subject identity intact, often does not accurately respond to the prompt conditions (\emph{e.g.}, see the initial row of the bowl instance). These qualitative results affirm that the OFT approach excels in both prompt controllability and fidelity of subject preservation. Additionally, OFT's performance is robust with respect to the number of finetuning iterations; even when the iteration count is increased, OFT remains effective—a property not shared by DreamBooth or LoRA. Additional qualitative outcomes can be found in Appendix~\ref{app:more_qualitative}. Further supporting evidence from human evaluation studies is presented in Appendix~\ref{app:human}, which corroborates OFT’s advantages.

\subsection{Controllable Generation}

\textbf{Settings}. We compare our approach to established baselines including ControlNet~\cite{zhang2023adding}, T2I-Adapter~\cite{mou2023t2i}, and LoRA~\cite{hulora2022}. The main evaluation focuses on three complex controllable generation tasks: converting Canny edge maps to images~(C2I) on the COCO dataset~\cite{lin2014microsoft}, transforming segmentation maps to images~(S2I) on ADE20K~\cite{zhou2017scene}, and mapping facial landmarks to face images~(L2F) using CelebA-HQ~\cite{karras2018progressive,xia2021tedigan}. For all experiments, each method is employed to finetune the Stable Diffusion (SD) v1.5 model on each dataset over 20 training epochs. Extended results are provided in Appendix~\ref{app:more_qualitative},\ref{app:more_control_task},\ref{app:failure}.

\textbf{Convergence}. We assess the convergence rates of ControlNet, T2I-Adapter, LoRA, and COFT on the L2F task, providing both quantitative and qualitative analyses. For quantitative assessment, we utilize the mean $\ell_2$ distance between the predicted face landmarks and the control face landmarks as our metric. Figure~\ref{fig:face_convergence} presents the face landmark errors recorded as the models are finetuned across varying epochs. The results reveal that both COFT and OFT converge at a much faster pace compared to alternatives. For instance, LoRA requires 20 epochs to reach the performance level that OFT attains by the 8-th epoch. It is important to mention that OFT and COFT possess a comparable number of trainable parameters to LoRA—substantially fewer than ControlNet—yet they reach convergence much more efficiently than other current approaches. Further, the rapid convergence of OFT is also corroborated by findings in Figure~\ref{fig:teaser}; in particular, the right-side example demonstrates that OFT exhibits far superior data efficiency relative to both ControlNet and LoRA, achieving robust convergence using only 5\% of the ADE20K dataset. For the qualitative evaluation, we emphasize comparisons among OFT, COFT, and ControlNet, as ControlNet most closely matches our landmark error rates. As seen in Figure~\ref{fig:face_convergence_qual}, both OFT and COFT demonstrate consistent convergence, with generated face poses progressively aligning with the control landmarks over time. By contrast, ControlNet displays a sudden onset of controllability post 8-th epoch, consistent with the abrupt convergence pattern discussed in \cite{zhang2023adding}. The stable convergence observed in OFT greatly facilitates its practical application, as it offers smoother and more predictable training behavior, thereby simplifying the process of tuning OFT’s hyperparameters.

\textbf{Quantitative comparison}. We propose a control consistency metric to quantitatively assess controllable generation. The central approach involves extracting the control signal from the generated output and comparing it to the original control input. In the case of the C2I task, we utilize IoU and F1 score as metrics. For S2I, we report mean IoU along with mean and overall accuracy. Regarding L2F, we measure the average $\ell_2$ distance between predicted and target control landmarks. Additional details about these consistency metrics are available in Appendix~\ref{app:settings}. Each method is evaluated using its optimal hyperparameter configuration. As presented in Table~\ref{table:control_gen}, OFT and COFT demonstrate significantly stronger and more precise control in comparison to alternative methods. Adapter-based strategies (such as T2I-Adapter and ControlNet) show slower convergence and inferior final outcomes. When comparing LoRA to ControlNet, LoRA exhibits improved performance on the S2I task, yet lags behind on C2I and L2F tasks. Overall, even after extensive hyperparameter tuning, LoRA's achievable performance remains limited. Conversely, our OFT framework appears to benefit from increasing the number of trainable parameters, with performance not yet reaching a plateau. It is important to note that our approach to quantitative evaluation in controllable generation constitutes a novel contribution, as it reliably measures the control capabilities of fine-tuned models (to the extent allowed by the pre-trained segmentation/detection model’s accuracy).

\textbf{Qualitative comparison}. We further conduct a qualitative evaluation of OFT and COFT against prevailing state-of-the-art techniques such as ControlNet, T2I-Adapter, and LoRA. As illustrated by randomly sampled results in Figure~\ref{fig:control_final}, both OFT and COFT are capable of producing images that are not only highly realistic and detailed, but also reflect precise controllability. In the S2I task, LoRA is unable to generate images that correspond to the provided segmentation maps, while ControlNet, OFT, and COFT are all able to appropriately guide image generation. When compared to ControlNet, OFT and COFT deliver images with superior fidelity, richer details, and more coherent geometric structures, while utilizing significantly fewer model parameters. For the C2I task, both OFT and COFT succeed in synthesizing semantically meaningful images even from rudimentary Canny edge inputs, a scenario where T2I-Adapter and LoRA underperform. In the context of the L2F task, our methods yield the highest accuracy in pose guidance, particularly under difficult facial poses. Across all three control tasks, OFT and COFT consistently outperform established baselines in terms of visual quality, validating the robust control capabilities of our OFT approach. To strengthen our qualitative assessment, we include a broader set of results for each control task in Appendix~\ref{app:control_exp}, and further showcase OFT's effectiveness in additional control settings (such as dense pose to human body, sketch-to-image, and depth-to-image) in Appendix~\ref{app:more_control_task}.

\section{Concluding Remarks and Open Problems}

Driven by insights into the pivotal role of angular relationships among neurons in shaping visual semantics, this work introduces a straightforward yet robust approach—orthogonal finetuning—for the manipulation of text-to-image diffusion models. Our focus is on two primary use cases: subject-driven generation and controllable generation. In comparison to current techniques, OFT offers enhanced controllability and stability during finetuning, requiring fewer tunable parameters. Notably, OFT maintains computational efficiency by avoiding extra inference cost, making it well-suited for real-world deployment.

This method also opens several noteworthy research questions. To ensure orthogonality, OFT utilizes Cayley parametrization, necessitating a matrix inversion, which can hinder its scalability. While our adoption of block diagonal parametrization mitigates this constraint, devising a differentiable and efficient mechanism to accelerate the matrix inversion remains unresolved. Additionally, OFT offers an intriguing avenue for compositionality: orthogonal matrices resulting from various OFT finetuning tasks can be successively multiplied while preserving orthogonality. It remains to be explored whether this ensemble of orthogonal matrices retains the accumulated knowledge from all associated tasks. Lastly, OFT's parameter efficiency is intrinsically linked to its block diagonal framework, which can introduce structural biases and reduce flexibility. Finding strategies to further improve parameter efficiency without sacrificing effectiveness or introducing new biases is an open problem worthy of further investigation.

\end{document}